% Nouveau EGA V (cote 25): archive cote 25, site page 120 (PDF page 121).
% typescript. Draft compared: none.
% AI-transcribed by Claude Opus 5.5 (claude-opus-5-5, adaptive thinking, effort high), Message Batches API, batch ega25b (msgbatch_01A7n5YuXubUPNK1zsepi8RG), 2026-11-13; checked against both supplied scan parts; see CHECK_DECISIONS.tsv.
\scanpage{25/120}
\mspage{\unsure{23|27}}

Z$_s$ est de codimension $\hand{\geq} 2$ dans X$_s$ ; \var{\xout{ou encore (comme dans les cas envisa}}{}{typewriter strike-out}{ega25_25_120-v0001}
$\hand{\unsure{\phi|\emptyset}}$ est alors l'ensemble des \var{\xout{familles}}{}{typewriter strike-out}{ega25_25_120-v0002} \var{\inserted[typed]{syst lin}}{syst lin}{insertion}{ega25_25_120-v0003} de diviseurs par.par P qui so\var{\inserted{\unsure{ie}}}{\unsure{ie}}{insertion}{ega25_25_120-v0004}nt
\uline{sans composante fixe}, \var{\xout{tandis que $\hand{\unsure{\psi|\emptyset}}$ est l'ensemble de}}{}{typewriter strike-out}{ega25_25_120-v0005} D'autre part, si
S=Spec(k), et si X est normal, alors \var{\inserted[typed]{pour}}{pour}{insertion}{ega25_25_120-v0006} \uline{toute} application rationelle
f:X$\hand{\to}$P, l'ensemble de définition U(f) vérifie codim ((X-U(f),X)$\hand{\geq}$2
(II 7), donc dans ce cas $\hand{\unsure{\psi|\emptyset}}$ est formé de l'ensemble de \uline{toutes} les
applications rationelles de X dans P.

La démonstration de a) a déjà été donnée. Pour prouver b), notons
que la formation de $i_{\hand{*}}(\underline{L}_U)$ commute avec toute extension plate de la
\var{\xout{pour le\unsure{x} il faut}}{}{typewriter strike-out}{ega25_25_120-v0007}
base S' \quad S (réf)\var{\inserted{du moins si $U\to X$ est quasi-compact, cas auquel on se ramène sans \unsure{difficultés}}}{du moins si $U\to X$ est quasi-compact, cas auquel on se ramène sans \unsure{difficultés}}{insertion}{ega25_25_120-v0008}, donc \var{\inserted[typed]{la cond. envisagée est inv. par}}{la cond. envisagée est inv. par}{insertion}{ega25_25_120-v0009} \var{\xout{il suffit de trouver un morphisme}}{}{typewriter strike-out}{ega25_25_120-v0010}
\var{\xout{S' $\quad$ S tel que la conclusion soit vraie après}}{}{typewriter strike-out}{ega25_25_120-v0011} changement de base \var{\inserted[typed]{fid plat qu cpc}}{fid plat qu cpc}{insertion}{ega25_25_120-v0012}. On
prendra S' = $\hand{\check{}}\mathrm{P}$ , et on note que l'hypothèse $i'_{\hand{*}}(\underline{L}'_U)$ inversible ne
change pas en remplaçant $\underline{L}'_U$ par $\underline{L}'_U\otimes_{S'}M'$, où M' est un Module inversible
sur S' (car $i'_{\hand{*}}(\underline{L}'_U\otimes_{S'}M')\hand{\simeq} i'_{\hand{*}}(\underline{L}'_U)\otimes_{S'}M'$ ). On \unsure{prèndra|prendra} M'= $\underline{O}_{\overset{\hand{\vee}}{P}}(1)$, de
sorte que la condition envisagée signifie aussi que $i'_{\hand{*}}(\underline{N}')$ est un Module
inversible, où $\underline{N}'= (f_U\times_S id_P)^{\hand{*}}(\underline{O}_{P\times\overset{\hand{\vee}}{P}}(1,1))$ . Or $\underline{O}$ (1,1) est précisément
le Module inversible défini par le diviseur canonique $\underline{H}$ sur $P\times_S\overset{\hand{\vee}}{P}$ , de sor%
te que $\underline{N}$ n'est autre que le Module inversible défini par \qquad D' =
$(f_U\times id_P)$ $(\underline{H})$. \var{\deleted{\var{\xout{\illegible{1}\unsure{D N P}\illegible{1}\unsure{Or}}}{}{typewriter strike-out}{ega25_25_120-v0014}\unsure{S}}}{}{deletion}{ega25_25_120-v0013} Lorsque f provient de D$\hand{\in}\hand{\unsure{\phi|\emptyset}}$ , \var{\xout{ce dernier}}{}{typewriter strike-out}{ega25_25_120-v0015}
D' n'est autre que D$\hand{|}$U$\times_S\overset{\hand{\vee}}{P}$ , \var{\xout{et il résulte du lemme}}{}{typewriter strike-out}{ega25_25_120-v0016} donc $\underline{N}' = \underline{N}\hand{|}U\times_S\overset{\hand{\vee}}{P}$ ,
où $\underline{N}$ est le Module inversible sur $X\times_S\hand{\check{}}P$ défini par D, et il résulte
du lemme \var{\inserted{3.2}}{3.2}{insertion}{ega25_25_120-v0017} plus haut \var{\xout{qu'on aura}}{}{typewriter strike-out}{ega25_25_120-v0018} appliqué à $P\times_S\hand{\check{}}P\hand{\to}\hand{\check{}}P$ qu'on aura
$i'_{\hand{*}}(\underline{N}')\hand{\simeq}\underline{N}$ , donc $i'_{\hand{*}}(\underline{N}')$ est inversible. Inversement, si cette condi-%
tion est remplie prouvons que f provient d'un\var{\xout{e}}{}{typewriter strike-out}{ega25_25_120-v0019} D$\hand{\in}\hand{\unsure{\phi|\emptyset}}$ , ou ce qui revient
ma\var{\xout{n}}{}{typewriter strike-out}{ega25_25_120-v0020}nifestement au même, que le diviseur D' se prolonge en un diviseur \var{\deleted{de}}{}{deletion}{ega25_25_120-v0021}
\var{\deleted{Cartier}}{}{deletion}{ega25_25_120-v0022} relatif par r apport \var{\xout{de}}{}{typewriter strike-out}{ega25_25_120-v0023} à $\hand{\check{}}P$ sur X$\times_S$P . Il \var{\deleted{suffit}}{}{deletion}{ega25_25_120-v0024} revient au même
de dire qu'il se prolonge en un diviseur \var{\deleted{de Cartier}}{}{deletion}{ega25_25_120-v0025} D sur X$\times_S$P , car \hand{D}.

\var{\margin{\var{\deleted{\unsure{Rq compliquée}}}{}{deletion}{ega25_25_120-v0027} \unsure{Il} faut \unsure{exprimer|expliquer} \illegible{1} $(f_U\times_S\mathrm{id}_P)^{-1}(H)$ \unsure{qui} est défini comme \unsure{diviseur} rel. à P \unsure{grâce} à la condition \illegible{3} \unsure{$f_U$} \unsure{se prolonge en} \unsure{un} diviseur relatif /P \unsure{à} $X\times_S P$ \unsure{/P} \var{\deleted{\illegible{2}}}{}{deletion}{ega25_25_120-v0028} \unsure{*}}}{\placedmargin{\var{\deleted{\unsure{Rq compliquée}}}{}{deletion}{ega25_25_120-v0027} \unsure{Il} faut \unsure{exprimer|expliquer} \illegible{1} $(f_U\times_S\mathrm{id}_P)^{-1}(H)$ \unsure{qui} est défini comme \unsure{diviseur} rel. à P \unsure{grâce} à la condition \illegible{3} \unsure{$f_U$} \unsure{se prolonge en} \unsure{un} diviseur relatif /P \unsure{à} $X\times_S P$ \unsure{/P} \var{\deleted{\illegible{2}}}{}{deletion}{ega25_25_120-v0028} \unsure{*}}}{marginal note placed}{ega25_25_120-v0026}

\mspage{120}
